\section{From the angular trace to the canonical source}
\label{sec:source}

\subsection{The scan, differential processing, and reflection}

The trace \eqref{eq:theta-trace} becomes a spatial profile only after an
external prescription. Set
\begin{equation}
 \beta(u)=\pi\ee^{4u},\qquad
 p(u)=\ee^uQ_1(\pi\ee^{4u}),\qquad
 \Phi(u)=\tfrac18(\partial_u^2-1)p(u).
 \label{eq:scan}
\end{equation}
For fixed $u$, differentiation of the Gaussian series is absolute and
locally uniform. With $x=\pi\ee^{4u}$ it gives
\begin{equation}
 \Phi(u)=\ee^u\sum_{m\ge1}
 \bigl(2(m^2x)^2-3m^2x\bigr)\ee^{-m^2x}.
 \label{eq:phi-series}
\end{equation}
The source summation index is written $m$ because the scalar source,
once prepared, is independent of the physical channel index $n$ of
the evolving field. Equating the scan coordinate $u$ to the physical
coordinate $v$ is also part of the preparation. The translation-invariant
Green kernel alone does not prescribe either operation.

The complete theta identity is what makes this source smooth and even.
For $R>0$, let
\begin{equation}
 \Theta(R)=\sum_{j\in\mathbb Z}\ee^{-\pi j^2R^2}
          =1+2Q_1(\pi R^2).
\end{equation}
The modular identity
$\Theta(R)=R^{-1}\Theta(R^{-1})$ is the zero-argument specialization
of \cite[20.7.32]{DLMFThetaModular}. Taking $R=\ee^{2u}$ yields
\begin{equation}
 p(u)-p(-u)=-\sinh u.
 \label{eq:reflection-p}
\end{equation}
The operator $\partial_u^2-1$ annihilates the right side, so
\begin{equation}
 \Phi(u)=\Phi(-u).
 \label{eq:phi-even}
\end{equation}
The series defining $p$ is smooth on every compact interval of the full
line; thus this is a smooth reflection, including the origin. Both the
zero-mode correction in $\Theta$ and the complete unit coefficients are
essential to this argument.

Every summand of \eqref{eq:phi-series} is positive when $u\ge0$, since
$m^2x\ge\pi>3/2$. Therefore $\Phi>0$ on the whole line. The first term
also gives, with $x=\pi\ee^{4|v|}$,
\begin{equation}
 \Phi(v)\ge\left(2-\frac3\pi\right)
 \ee^{|v|}x^2\ee^{-x}>0.
 \label{eq:phi-lower}
\end{equation}
This positivity concerns the source. It will not determine the sign of
its oscillatory Fourier observation.

\subsection{A hierarchy of complete-source bounds}

We next establish a single collection of majorants that will control the
preparation, evolution, and readout. Define polynomials
\begin{equation}
 P_0(x)=2x^2-3x,\qquad
 P_{j+1}(x)=(1-4x)P_j(x)+4xP_j'(x).
 \label{eq:source-poly}
\end{equation}
Then $\deg P_j=j+2$, and differentiating on $u\ge0$ gives
\begin{equation}
 \Phi^{(j)}(u)=\ee^u\sum_{m\ge1}
                 P_j(m^2x)\ee^{-m^2x}.
 \label{eq:phi-derivatives}
\end{equation}
Let $a_j$ be the sum of the absolute coefficients of $P_j$, and put
\begin{equation}
 C_d(N)=\sum_{m>N}m^{2d}\ee^{-\pi(m^2-1)},\qquad
 B_j=a_jC_{j+2}(0).
 \label{eq:source-constants}
\end{equation}
Every $C_d(N)$ is finite and tends to zero with $N$. Since $x\ge\pi>1$,
\begin{equation}
 \ee^u\sum_{m>N}|P_j(m^2x)|\ee^{-m^2x}
 \le a_j C_{j+2}(N)\ee^u x^{j+2}\ee^{-x}.
 \label{eq:source-tail}
\end{equation}
Indeed, $|P_j(m^2x)|\le a_jm^{2(j+2)}x^{j+2}$ and
$\ee^{-(m^2-1)x}\le\ee^{-\pi(m^2-1)}$. Thus the estimate controls
the sum of absolute values, not merely cancellation in a signed tail.
Reflection gives the global bound
\begin{equation}
 |\Phi^{(j)}(v)|\le B_j\ee^{|v|}x^{j+2}\ee^{-x},
 \qquad x=\pi\ee^{4|v|}.
 \label{eq:phi-global}
\end{equation}

For nonnegative integers $r,d$ and finite $M,b\ge0$, define
\begin{equation}
 J_{r,d}(M,b)=\int_\R |v|^r
 \ee^{Mv^2+(b+1)|v|}x^d\ee^{-x}\dd v.
 \label{eq:J-definition}
\end{equation}
These integrals are finite. More explicitly, the substitution
$x=\pi\ee^{4|v|}$ gives
\begin{equation}
 J_{r,d}(M,b)=\frac{2}{4^{r+1}}\int_\pi^\infty
 [\log(x/\pi)]^r(x/\pi)^{(b+1)/4}x^{d-1}
 \exp\!\left(\frac{M}{16}\log^2(x/\pi)-x\right)\dd x.
 \label{eq:J-integral}
\end{equation}
The term $-x$ dominates the logarithmic square and all polynomial factors.
The remaining integral beyond a radial cutoff consequently tends to
zero. This provides a complete radial tail bound at every fixed order.

Write
\begin{equation}
 F_\lambda(v)=\ee^{\lambda v^2}\Phi(v),\qquad
 f_{a,\lambda}(v)=v^{2a}F_\lambda(v)=\partial_\lambda^aF_\lambda(v).
 \label{eq:deformed-source}
\end{equation}
Define $R_0(\lambda,v)=1$ and
$R_{j+1}=\partial_vR_j+2\lambda vR_j$, so that
$\partial_v^j\ee^{\lambda v^2}=\ee^{\lambda v^2}R_j(\lambda,v)$.
For each fixed $a,h$ and $M$, finite Leibniz expansion gives
\begin{equation}
 |f_{a,\lambda}^{(h)}(v)|
 \le C_{a,h,M}(1+|v|)^{2a+h}
 \ee^{Mv^2+|v|}x^{h+2}\ee^{-x},\qquad |\lambda|\le M.
 \label{eq:f-majorant}
\end{equation}
For specificity, let $r_j(M)$ be the sum of the absolute coefficients
of $R_j$ after replacing every power $\lambda^t$ by $M^t$. One permissible
constant is
\begin{equation}
 C_{a,h,M}=\sum_{\substack{i+j+\ell=h\\i\le2a}}
 \frac{h!}{i!j!\ell!}\frac{(2a)!}{(2a-i)!}\,r_j(M)B_\ell.
 \label{eq:f-constant}
\end{equation}
Here $|R_j(\lambda,v)|\le r_j(M)(1+|v|)^j$, and
$x^{\ell+2}\le x^{h+2}$; these inequalities prove
\eqref{eq:f-majorant} directly.

It is useful to collect the weighted $L^1$ bounds as
\begin{align}
 w_{r,b}(v)&=(1+|v|)^r\ee^{b|v|},\\
 \mathcal N_{a,h;r,b}(M)
 &=\sup_{|\lambda|\le M}\norm{w_{r,b}f_{a,\lambda}^{(h)}}_1
 \nonumber\\
 &\le C_{a,h,M}\sum_{\ell=0}^{2a+h+r}
       \binom{2a+h+r}{\ell}J_{\ell,h+2}(M,b)<\infty.
 \label{eq:N-bound}
\end{align}
Squaring \eqref{eq:f-majorant} proves all finite $L^2$ derivative
bounds. Multiplying it by any polynomial and taking a supremum proves
every fixed Schwartz seminorm. The same majorants, with one further
factor $v^2$, prove continuity and differentiation in $\lambda$ in
each of these spaces. In particular, the image of a compact $\lambda$
interval is compact in every one of the specified norms.

Two distinctions about completeness will be used later. First, a
reflected finite source prefix $\Phi_N$ obeys the value-and-moment bound
\begin{equation}
 \norm{w_{r,b}v^{2a}\ee^{\lambda v^2}(\Phi-\Phi_N)}_1
 \le a_0C_2(N)\sum_{\ell=0}^r\binom r\ell
                      J_{2a+\ell,2}(M,b)\longrightarrow0.
 \label{eq:prefix-value-tail}
\end{equation}
Such a prefix generally has derivative jumps at the reflection point.
It is not a global smooth graph-core approximation; the operator cores
of Proposition~\ref{prop:complete-realization} supply those approximations.
Second, integration by parts yields, for any finite $d,l\ge0$,
\begin{equation}
 |\partial_k^d\widehat f_{a,\lambda}(k)|
 \le\frac{(1+k^2)^{-l}}{\sqrt{2\pi}}
 \norm{(1-\partial_v^2)^l[(-\ii v)^df_{a,\lambda}]}_1.
 \label{eq:fourier-tail}
\end{equation}
The right side is uniformly finite by \eqref{eq:N-bound}. Choosing $l$
sufficiently large controls every fixed polynomially weighted radial
Fourier tail in $L^1$ or $L^2$.

\subsection{The Newman transform and its absolute normalization}

Evenness converts the cosine transform into
\begin{equation}
 H_\lambda(z)=\frac12\int_\R F_\lambda(v)\ee^{-\ii zv}\dd v,
 \qquad \widehat F_\lambda(k)=\sqrt{\frac2\pi}\,H_\lambda(k).
 \label{eq:H-fourier}
\end{equation}
On $|\lambda|\le M$, $|\Im z|\le b$, the source estimates justify
every fixed finite mixed derivative:
\begin{equation}
 \partial_\lambda^a\partial_z^d H_\lambda(z)
 =\frac12\int_\R v^{2a}(-\ii v)^dF_\lambda(v)
                   \ee^{-\ii zv}\dd v,
 \qquad
 |\partial_\lambda^a\partial_z^d H_\lambda(z)|
 \le\frac{B_0}{2}J_{2a+d,2}(M,b).
 \label{eq:H-derivatives}
\end{equation}
The same domination on compact complex $\lambda$ sets shows joint
entire dependence on $(\lambda,z)$. In particular,
\begin{equation}
 \partial_\lambda H_\lambda=-\partial_z^2H_\lambda.
 \label{eq:newman-heat}
\end{equation}
The minus sign comes from $(-\ii v)^2=-v^2$. The parameter $\lambda$
reweights the source; it is not the cylinder heat parameter.

For completeness, the factor $1/8$ in the zeta normalization can be
derived directly, rather than left as an adjustable gain. By evenness we
may use $\ee^{\ii zv}$ in \eqref{eq:H-fourier}. Put
$s=(1+\ii z)/2$ and $t=\ee^{4v}$. For $\Re s>1$, absolute summability
allows termwise integration over the full line and gives
\begin{align}
 H_0(z)
 &=\frac18\sum_{m\ge1}\int_0^\infty
 t^{s/2-1}\bigl[2(\pi m^2t)^2-3\pi m^2t\bigr]
                         \ee^{-\pi m^2t}\dd t
 \nonumber\\
 &=\frac{2\Gamma(s/2+2)-3\Gamma(s/2+1)}8
                         \pi^{-s/2}\zeta(s)
 \nonumber\\
 &=\frac{s(s-1)}{16}\Gamma(s/2)\pi^{-s/2}\zeta(s)
  =\frac18\xi(s).
 \label{eq:xi-normalization}
\end{align}
Absolute integrals of the individual terms are bounded by a constant
times $m^{-\Re s}$ in this half-plane. The already established entire
integral supplies the continuation of the identity to all $z$.
This agrees with \cite[Introduction, eqs.~(1)--(4)]{RodgersTao2020}.
For real $\lambda$, $H_\lambda(0)>0$ by source positivity, so none of
these entire functions is identically zero.
